\section{Auditoría de fuentes: esto no es una tabla universal de recetas de Alignment}

Lo más valioso de esta clase no son unos cuantos hiperparámetros fijos, sino una forma de organizar los experimentos: primero se descompone el «helpful chatbot» en demonstration data, preference data, training objective y evaluation protocol; después se cambian, uno por uno, la distribución de los datos y el evaluador, y se observa si las conclusiones se mantienen estables. La versión anterior comprimía estos niveles en una lista de nombres (SFT, reward model, PPO/DPO) e incluía una tasa de aprendizaje y un batch size uniformes que la clase no respalda; esta reescritura vuelve a la evidencia de la clase del 31 de octubre de 2023 y coloca las salvedades del ponente junto a cada conclusión.

\subsection{Materiales oficiales e identificación de versiones}

Las fuentes principales de esta sesión son la página del curso Stanford CS25 V3, la grabación oficial de Stanford Online \texttt{mcep6W8oB1I}, los subtítulos manuales \texttt{en-US} y el deck público de la ponente. En el sitio de la ponente existen dos versiones, de 67 y de 71 páginas; la grabación proyecta con claridad ``Recipe 2'' y ``Zephyr-7B'' entre 00:26:30--00:26:40, por lo que se eligió el \texttt{transformers\_united.pdf} de 71 páginas, que contiene esas diapositivas, y no la variante de 67 páginas, a la que le faltan las páginas de distillation y de UN advisory.

\lecturefigure{slide-01.jpg}{Portada del deck oficial: Recipes for Training Helpful Chatbots}{Diapositivas oficiales p.1}

\begin{knowledgebox}{Lectura de la figura: Recipe designa una descomposición de ingeniería, no un ajuste misterioso}
El plural ``Recipes'' del título ya indica que no existe un procedimiento único. Al leer cada grupo de experimentos conviene hacerse cuatro preguntas: quién escribió las muestras, qué objetivo se optimiza, qué evaluador se usa y si la conclusión se reproduce con otros evaluadores. Si solo se registran los nombres de los modelos y las tablas de clasificación, se pierde el verdadero método de investigación de esta sesión.
\end{knowledgebox}

\lecturefigure{slide-02.jpg}{El equipo H4, sus objetivos, ingredients y procedure}{Diapositivas oficiales p.2}

\begin{importantbox}{Las cuatro H de H4}
H4 se refiere a \term{Helpful, Harmless, Honest, Huggy}. Aunque la sesión parte de estos cuatro objetivos, estudia principalmente la helpfulness y los SFT data; la harmlessness, la honesty y una safety evaluation completa aparecen solo en los límites de preference/red-teaming, por lo que esta sesión no debe equipararse con una solución de seguridad integral.
\end{importantbox}

\teachervoice{Desde el inicio, Rajani explica que el equipo quiere reproducir la ``secret sauce'' del alignment sobre un open-access pretrained model, pero la mayor parte de la clase se centra en el primer paso: cómo convertir el modelo base en un helpful chatbot capaz de completar las tareas del usuario. Esta delimitación del alcance fija los límites de interpretación de todos los resultados posteriores.}

\lecturefigure{slide-03.jpg}{Ruta de la clase: SFT, datos de RLHF, distillation, experimentos y sesgos del evaluador}{Diapositivas oficiales p.3}

\begin{knowledgebox}{Lectura de la figura: la sesión tiene dos recipes y una línea de auditoría}
La primera recipe consiste en human-curated demonstrations/preferences; la segunda, en los AI-generated data y el AI feedback de Zephyr; ambas rutas desembocan finalmente en la evaluation. Los GPT-4 evaluator quirks del extremo derecho no son un apéndice, sino una auditoría inversa de todas las conclusiones anteriores basadas en tablas de clasificación.
\end{knowledgebox}

\subsection{Límites de la evidencia y estrategia de cobertura}

De las 71 páginas del deck, 66 contienen información didáctica; se omiten a propósito las páginas separadoras 6, 32, 48 y 59 y la página final de Thanks (71), y se conservan todos los demás progressive highlights. Las tablas de clasificación, los costos y las capacidades de los modelos citados en el texto son lecture-time snapshots: por ejemplo, ``Zephyr beats ChatGPT'' solo significa que obtuvo un resultado específico con los protocolos de AlpacaEval/MT-Bench mostrados en ese momento; no demuestra que sea más seguro, más veraz ni superior en todas las tareas.

\begin{warningbox}{No conviertas un snapshot de 2023 en una timeless law}
Las tablas de clasificación cambian con el prompt template, el judge model, el sampling, la filtración de datos y los modelos participantes. La conclusión más duradera de esta sesión no es la clasificación de ningún modelo, sino esta: los experimentos con datos deben leerse junto con un evaluator audit; de lo contrario, una respuesta más larga o más parecida a la distribución de entrenamiento del judge puede tomarse erróneamente por más helpful.
\end{warningbox}

\subsection{Resumen de la sección}

Esta sesión es un experimento conjunto de data--objective--evaluator, no una tabla general de ajuste de hiperparámetros. La forma correcta de leerla es distinguir siempre entre los hechos de la clase, los juicios de ingeniería del ponente, las métricas sustitutas de las tablas de clasificación y las explicaciones de mecanismos que añadimos con base en los artículos originales.
